\documentclass[12pt,a4paper]{article}
\usepackage[polish]{babel}
\usepackage[T1]{fontenc}
\usepackage[utf8x]{inputenc}
\usepackage{hyperref}
\usepackage{url}
\usepackage[]{algorithm2e}
\usepackage{listings}
\usepackage{amsmath}
\usepackage{graphicx}

\usepackage{color}
\usepackage{listings}

\lstloadlanguages{% Check Dokumentation for further languages ...
	C,
	C++,
	csh,
	Java
}

\definecolor{red}{rgb}{0.6,0,0} % for strings
\definecolor{blue}{rgb}{0,0,0.6}
\definecolor{green}{rgb}{0,0.8,0}
\definecolor{cyan}{rgb}{0.0,0.6,0.6}

\lstset{
	language=Python,
	basicstyle=\footnotesize\ttfamily,
	numbers=left,
	numberstyle=\tiny,
	numbersep=5pt,
	tabsize=2,
	extendedchars=true,
	breaklines=true,
	frame=b,
	stringstyle=\color{blue}\ttfamily,
	showspaces=false,
	showtabs=false,
	xleftmargin=17pt,
	framexleftmargin=17pt,
	framexrightmargin=5pt,
	framexbottommargin=4pt,
	commentstyle=\color{green},
	morecomment=[l]{\#},
	showstringspaces=false,
	keywordstyle=\color{cyan},
	identifierstyle=\color{red},
}
\usepackage{caption}
\DeclareCaptionFont{white}{\color{white}}
\DeclareCaptionFormat{listing}{\colorbox{blue}{\parbox{\textwidth}{\hspace{15pt}#1#2#3}}}
\captionsetup[lstlisting]{format=listing,labelfont=white,textfont=white, singlelinecheck=false, margin=0pt, font={bf,footnotesize}}


\addtolength{\hoffset}{-1.5cm}
\addtolength{\marginparwidth}{-1.5cm}
\addtolength{\textwidth}{3cm}
\addtolength{\voffset}{-1cm}
\addtolength{\textheight}{2.5cm}
\setlength{\topmargin}{0cm}
\setlength{\headheight}{0cm}

\begin{document}
	
	\title{Systemy Sztucznej Inteligencji\\
	\small{Dokumentacja projektu -- Klasyfikacja przedziału cenowego rosyjskich mieszkań}}
	\author{Kamil Skop, Michał Kosiec, Szymon Sobelski - gr. 8, ?, ?}
	\date{\today}

	\maketitle
	\tableofcontents
	\newpage

	%--------------------------------------------------------------
	\section{Część I}
	%--------------------------------------------------------------

	\subsection{Opis programu}

	Celem projektu jest klasyfikacja przedzialu cenowego rosyjskich mieszkan na podstawie ich cech
	(powierzchnia, liczba pokoi, pietro, lokalizacja geograficzna, typ budynku itp.).
	Dane wejsciowe stanowi plik CSV zawierajacy 200\,000 rekordow.
	Cena kazego mieszkania zostaje zakwalifikowana do jednego z $N = 5$ rowno licznych przedzialow cenowych
	(klas \texttt{C1}--\texttt{C5}) wyznaczanych dynamicznie kwantylami ze zbioru treningowego.
	
	Program losowo, z użyciem stałego \texttt{seed = 42} dzieli wczytane dane na zbiór treningowy oraz testowy w proporcji \texttt{80:20} po uprzednim oczyszczeniu go z outlierów na podstawie reguły 3-sigma na cechach ciągłych (gdy odległość dowolnej cechy od średniej przekracza 3 odchylenia standardowe, usuwany jest cały rekord).

	Program implementuje od zera, bez uzycia zewnetrznych bibliotek uczenia maszynowego,
	trzy klasyfikatory:
	\begin{itemize}
		\item \textbf{KNN} (k najblizszych sasiadow),
		\item \textbf{Naiwny Bayes} (mieszany: Gauss + Bernoulli + wielomianowy),
		\item \textbf{Drzewo decyzyjne} (kryterium Gini, ograniczenie glebokosci i minimalnego liscia).
	\end{itemize}
	
	W sprawozdaniu przeprowadzono również analizę rezultatów każdego z wytrenowanych klasyfikatorów.
	
	\newpage
	\subsection{Instrukcja obsługi}
	\textbf{Wymagania:}
	\begin{itemize}
		\item Python 3.12+
		\item Biblioteka \texttt{matplotlib}
		\item Obecność przetworzonej bazy danych z Kaggle
	\end{itemize}
	
	\textbf{Uruchomienie:}
	\begin{itemize}
		\item Upewnić się, że w katalogu projektu znajduje się plik \texttt{filtered\_v3.csv}. Jeśli ten nie występuje, należy pobrać zbiór danych ze strony Kaggle, a następnie wygenerować go za pośrednictwem skryptu, umieszczonego w katalogu \texttt{datalab/}.
		\item Uruchomić skrypt za pomocą polecenia \texttt{python main.py}.
	\end{itemize}
	
	\textbf{Uwaga:} Przetestowanie KNN może wymagać zmniejszenia parametru \texttt{EXPECTED\_ROWS} w kodzie źródłowym do np. 10 000, ponieważ algorytm ten jest bardzo niewydajny i zwyczajne nie radzi sobie z większymi ilościami rekordów.
	
	\subsection{Opis obsługi}
	Program wczytuje plik \texttt{filtered\_v3.csv} z biezacego katalogu,
	usuwa outliery i prezentuje menu tekstowe. Należy wybrać interesującą nas opcję, wprowadzając ją z klawiatury i wciskając enter:
	\begin{figure}[htbp]
		\centering
		\includegraphics[width=0.79\textwidth]{img/main\_menu.png}
		\caption{Zrzut ekranu menu głównego.}
		\label{fig:menu_glowne}
	\end{figure}
	\begin{enumerate}
		\item Analiza danych (wyświetla analizę danych oraz generuje wykresy do katalogów \texttt{output\_raw/} i \texttt{output\_clean/}).
		\item Trening modelu KNN.
		\item Trening modelu Naiwny Bayes.
		\item Trening modelu Drzewo decyzyjne.
		\item Analiza wyników aktualnie wytrenowanego modelu (macierz pomyłek, dokładność itd.).
		\item Wyjście.
	\end{enumerate}

	\newpage

	%--------------------------------------------------------------
	\section{Część II}
	%--------------------------------------------------------------

	\subsection{Opis działania}
	
	\subsubsection{Przetwarzanie wstępne i binning cen}
	
	Cena mieszkania jest wartością ciągłą. Aby problem stał się klasyfikacją,
	ceny dzielone są na $N$ równo licznych przedziałów (kwantylowych):
	\[
	q_i = \text{kwantyl}\!\left(\frac{i}{N}\right), \quad i = 1,\ldots, N-1.
	\]
	Dla $N = 5$ otrzymujemy 5 klas \texttt{C1}--\texttt{C5}.
	Podział wyznaczany jest raz z całego zbioru (po usunięciu outlierów),
	co gwarantuje równo liczne klasy i eliminuje wpływ wartości ekstremalnych na granice.
	
	\subsubsection*{Usuwanie outlierów -- reguła 3-sigma}
	
	Dla każdej cechy liczbowej $x_j$ wyznaczana jest średnia $\mu_j$ i odchylenie standardowe $\sigma_j$.
	Rekord uznaje się za outlier, jeżeli dla dowolnej cechy $j$:
	\[
	|x_j - \mu_j| > 3\sigma_j.
	\]
	
	\subsubsection{Algorytm KNN}
	
	Klasyfikator k najbliższych sąsiadów ($k = 7$) przydziela klasę na podstawie
	ważonego głosowania sąsiadów. Dla pary przykładów $a$, $b$ stosuje się metrykę hybrydową:
	
	\[
	d(a, b) = \underbrace{\sqrt{\sum_{j} \left(\frac{a_j - b_j}{s_j}\right)^2}}_{\text{część ciągła}}
	+ d_{\text{region}} + d_{\text{budynek}} + d_{\text{nowe}}
	\]
	gdzie $s_j$ to odchylenie standardowe $j$-tej cechy ciągłej (normalizacja),
	a dodatkowe składniki binarne karzą za różny region ($+1{,}0$), typ budynku ($+0{,}5$)
	lub różny status nowego budownictwa ($+0{,}4$).
	
	Waga sąsiada odwrotnie proporcjonalna do odległości:
	\[
	w_i = \frac{1}{d_i + \varepsilon}, \quad \varepsilon = 10^{-9}.
	\]
	
	\newpage
	\subsubsection{Naiwny Bayes}
	
	Klasyczny Naiwny Bayes zakłada, że wszystkie cechy są warunkowo niezależne
	i mają rozkład normalny. Diagnoza wizualna rozkładów cech wykazała, że założenie
	o gaussowskości jest niespełnione dla większości zmiennych:
	
	\begin{itemize}
		\item \textbf{Lokalizacja geograficzna}: rozkład wielomodalny (skupiska miejskie).
		\item \textbf{Miesiąc}: rozkład cykliczny.
		\item \textbf{Liczba pokoi / pięter}: wartości dyskretne ze skokami.
	\end{itemize}
	
	W odpowiedzi wdrożono \emph{mieszany Naiwny Bayes} (Mixed Naive Bayes)
	z trzema niezależnymi torami prawdopodobieństwa:
	
	\begin{enumerate}
		\item \textbf{Gaussowski} -- dla $\log(\text{area}+\varepsilon)$ i $\log(\text{kitchen\_area}+\varepsilon)$:
		\[
		P(x_j \mid c) = \frac{1}{\sqrt{2\pi\sigma_{jc}^2}}
		\exp\!\left(-\frac{(x_j - \mu_{jc})^2}{2\sigma_{jc}^2}\right).
		\]
		\item \textbf{Bernoulliego} -- dla cechy binarnej \texttt{new\_building}
		(wygładzanie Laplace'a, $\alpha = 1$):
		\[
		P(x_j \mid c) = p_{jc}^{x_j}(1-p_{jc})^{1-x_j}, \quad
		p_{jc} = \frac{n_{1jc} + 1}{n_c + 2}.
		\]
		\item \textbf{Wielomianowy} -- dla pozostałych cech dyskretnych i kategorycznych
		(rok, miesiąc, piętro, liczba pięter, pokoje, region, typ budynku, hasz geo):
		\[
		P(x_j = v \mid c) = \frac{n_{vjc} + \alpha}{n_c + \alpha \cdot |V_j|}.
		\]
	\end{enumerate}
	
	Spatial hashing: współrzędne geograficzne zaokrąglane są do 0{,}1 stopnia
	($\approx 11$ km), co zastępuje kosztowne klasteryzowanie i pozwala traktować
	lokalizacje jako kategorie.
	
	Klasyfikacja: wybierana jest klasa z maksymalnym logarytmem a posteriori:
	\[
	\hat{c} = \arg\max_c \left[
	\log P(c) + \sum_j \log P(x_j \mid c)
	\right].
	\]
	
	\subsubsection{Drzewo decyzyjne}
	
	Drzewo budowane jest metodą zachłanną (CART). Kryterium podziału to nieczystość Gini:
	\[
	\text{Gini}(S) = 1 - \sum_{c} \left(\frac{|S_c|}{|S|}\right)^2.
	\]
	Najlepszy podział minimalizuje ważoną sumę nieczystości potomków:
	\[
	\text{Gain} = \text{Gini}(S) - \frac{|S_L|}{|S|}\text{Gini}(S_L) - \frac{|S_R|}{|S|}\text{Gini}(S_R).
	\]
	Ograniczenia: maksymalna głębokość $= 13$, minimalna liczba próbek w liściu $= 10$.
	Trening odbywa się na losowej próbce 5\,000 rekordów (seed = 42).

	\subsection{Algorytm}

	Ponizej przedstawiono pseudokod glownego potoku przetwarzania danych.

	\begin{algorithm}[H]
		\KwData{Plik CSV z danymi mieszkan, liczba klas $N = 5$, proporcja treningu $r = 0{,}8$}
		\KwResult{Wytrenowany klasyfikator, wyniki na zbiorze testowym}
		Wczytaj wszystkie rekordy z pliku CSV\;
		Wyznacz granice klas cenowych kwantylami z cen wszystkich rekordow\;
		Przypisz kazdy rekord do klasy \texttt{C1}--\texttt{C5}\;
		Oblicz srednia i odch. stand. kazdej cechy numerycznej\;
		\ForEach{rekord}{
			\If{dowolna cecha odbiega o $> 3\sigma$}{
				Oznacz jako outlier\;
			}
		}
		Usun outliery ze zbioru\;
		Przetasuj zbior losowo (seed=42)\;
		Podziel na zbior treningowy ($80\%$) i testowy ($20\%$)\;
		\textbf{Wybierz model} (KNN / Naiwny Bayes / Drzewo decyzyjne)\;
		Trenuj wybrany model na zbiorze treningowym\;
		\ForEach{przyklad ze zbioru testowego}{
			Przewidz klase cenowa\;
			Porownaj z etykieta rzeczywista\;
		}
		Wyswietl macierz pomylek i dokladnosc\;
		\caption{Główny potok algorytmu programu}
	\end{algorithm}

	\newpage
	\subsection{Baza danych}
	Dane wejściowe: plik \texttt{filtered\_v3.csv} (ok. 15 MB, 200\,000 rekordów). Zbiór został uzyskany poprzez przetasowanie i sztuczne zmniejszenie zbioru bazowego, pochodzącego z Kaggle, link: \href{https://www.kaggle.com/datasets/mrdaniilak/russia-real-estate-20182021}{https://www.kaggle.com/datasets/mrdaniilak/russia-real-estate-20182021} za pomocą skryptu z katalogu \texttt{datalab/}. Zawiera on około 5,5 miliona rekordów ogłoszeń sprzedaży mieszkań w Rosji w latach 2018 - 2021, redukowanych i tasowanych do 200 000 wpisów. Jako, że w plikach projektu nie załączono bazy, należy wygenerować ją samodzielnie z użyciem wspomnianego wyżej skryptu oraz pliku pobranego z linku.
	
	\subsection{Struktura rekordu}
	\begin{center}
	\begin{tabular}{|l|l|l|}
		\hline
		\textbf{Kolumna} & \textbf{Typ} & \textbf{Opis} \\
		\hline
		price & int & Cena mieszkania (RUB) \\
		date & date & Data wystawienia ogłoszenia (YYYY-MM-DD) \\
		time & time & Czas wystawienia ogłoszenia \\
		geo\_lat & float & Szerokość geograficzna \\
		geo\_lon & float & Dlugość geograficzna \\
		region & int & ID regionu Rosji \\
		building\_type & int & Typ budynku (więcej informacji dostępnych na Kaggle) \\
		level & int & Piętro mieszkania \\
		levels & int & Liczba pięter w budynku \\
		rooms & int & Liczba pokoi (-1 = kawalerka) \\
		area & float & Powierzchnia całkowita (m$^2$) \\
		kitchen\_area & float & Powierzchnia kuchni (m$^2$) \\
		object\_type & int & 1 = rynek wtórny, 11 = nowe budownictwo \\
		\hline
	\end{tabular}
	\end{center}

	\newpage
	\subsection{Implementacja}
	Projekt podzielony jest na nastepujące pliki:
	\begin{itemize}
		\item \texttt{russian\_flat.py} -- klasy danych (\texttt{RussianFlatCsv}, \texttt{RussianFlatProps}, \texttt{RussianFlat}); wczytywanie i etykietowanie rekordów.
		\item \texttt{preprocessor.py} -- usuwanie outlierów (reguła 3-sigma) i losowy podział na zbiory.
		\item \texttt{base\_model.py} -- abstrakcyjna klasa bazowa \texttt{BaseModel}.
		\item \texttt{knn.py} -- klasyfikator KNN z hybrydową metryką i ważoną głosowaniem.
		\item \texttt{naive\_bayes.py} -- mieszany Naiwny Bayes (Gauss + Bernoulli + wielomianowy).
		\item \texttt{decision\_tree.py} -- drzewo decyzyjne CART z kryterium Gini.
		\item \texttt{data\_analysis.py} -- analiza statystyczna, macierz korelacji Pearsona, wykresy.
		\item \texttt{main.py} -- główna petla menu, orkiestracja całego potoku.
	\end{itemize}
	
	\subsubsection*{Wybrane fragmenty kodu}

	\textbf{Przykład:} Wyznaczanie skal do normalizacji w KNN.
	\begin{lstlisting}[caption=Obliczanie skal normalizacyjnych (knn.py)]
def _compute_continuous_scales(self, training_data):
    examples = [_ for _ in training_data]
    columns = list(zip(*(
        self._extract_features(flat.data)[0]
        for flat in examples
    )))
    scales = []
    for column in columns:
        stdev = statistics.stdev(column) if len(column) > 1 else 0.0
        if stdev <= 0.0:
            stdev = 1.0
        scales.append(stdev)
    return tuple(scales)
	\end{lstlisting}

	\textbf{Przykład:} Trening Naiwnego Bayesa -- wygładzanie wariancji dla toru gaussowskiego.
	\begin{lstlisting}[caption=Wygładzanie wariancji (naive\_bayes.py)]
global_variances = []
global_columns = list(zip(*all_continuous_features))
for col in global_columns:
    mean = sum(col) / len(col)
    var = sum((x - mean)**2 for x in col) / (len(col) - 1)
    global_variances.append(var)
epsilons = [max(1e-4 * var, 1e-9) for var in global_variances]
	\end{lstlisting}

	\subsection{Eksperymenty}
	
	Poniżej przedstawiono wyniki przeprowadzonych eksperymentów, obejmujących trening przygotowanych klasyfikatorów oraz ich ocenę na zbiorze testowym.
	
	\subsubsection*{KNN}
	
	\subsubsection*{Naiwny Bayes}
	
	\subsubsection*{Drzewo decyzyjne}
	
	\subsection*{Porównanie}

	\subsection{Wnioski}

	Na podstawie wygenerowanych histogramów z nakładaną krzywą Gaussa stwierdzono:
	\begin{itemize}
		\item Cechy \texttt{area} i \texttt{kitchen\_area} -- po logarytmowaniu przybliżają rozkład normalny; traktowane gaussowsko.
		\item Cechy \texttt{geo\_lat}, \texttt{geo\_lon} -- wyraźny rozkład wielomodalny (skupiska miast); traktowane jako kategorie przez spatial hashing.
		\item Cechy \texttt{publish\_month} -- rozkład cykliczny; traktowane jako kategorie.
		\item Cechy \texttt{level}, \texttt{levels}, \texttt{rooms} -- wartości dyskretne; traktowane wielomianowo.
	\end{itemize}
	
	Na podstawie wyników eksperymentów stwierdzono:
	\begin{itemize}
		\item W praktyce, najlepszym z wytrenowanych klasyfikatorów okazał się Naiwny Bayes ze względu na swoją wysoką wydajność oraz skuteczność, a także brak problemów z treningiem, biorącym pod uwagę cały zbiór danych wejściowych.
		\item KNN ze względu na swoją bardzo niską wydajność, pomimo osiągnięcia najwyższej dokładności, w praktycznych zastosowaniach byłby bezuzyteczny.
	\end{itemize}

	\newpage

	%--------------------------------------------------------------
	\section{Pełen kod aplikacji}
	%--------------------------------------------------------------

	\lstset{language=Python, basicstyle=\scriptsize\ttfamily, breaklines=true,
		numbers=left, numberstyle=\tiny, frame=b,
		commentstyle=\color{green}, keywordstyle=\color{cyan},
		identifierstyle=\color{red}, stringstyle=\color{blue}\ttfamily}

	\subsubsection*{main.py}
	\begin{lstlisting}
# Trzeba tu wrzucic wszystkie nasze pliki Pythona
print("Hello World!")
	\end{lstlisting}

	\subsubsection*{knn.py}
	\begin{lstlisting}
print("Hello World!")
	\end{lstlisting}

\end{document}
